% =========================================================
% Section: Boolean Rings — Definitions and First Properties
% Chapter: Boolean Rings and Boolean Algebras — Volume I
% Source: Halmos & Givant, Introduction to Boolean Algebras, §1
% =========================================================

% ---------------------------------------------------------
% The Two-Element Boolean Ring 2
% ---------------------------------------------------------

The simplest non-degenerate Boolean ring is the two-element ring
$\mathbf{2} = \{0, 1\}$.
Its addition and multiplication are given by the following tables.

\begin{center}
\begin{tabular}{c|cc}
$+$ & $0$ & $1$ \\ \hline
$0$ & $0$ & $1$ \\
$1$ & $1$ & $0$
\end{tabular}
\qquad\qquad
\begin{tabular}{c|cc}
$\cdot$ & $0$ & $1$ \\ \hline
$0$ & $0$ & $0$ \\
$1$ & $0$ & $1$
\end{tabular}
\end{center}

\noindent
Addition in $\mathbf{2}$ is addition modulo $2$: $1 + 1 = 0$.
Multiplication in $\mathbf{2}$ is ordinary multiplication of
integers restricted to $\{0, 1\}$.
Every element of $\mathbf{2}$ satisfies $p + p = 0$ and
$p \cdot p = p$, making $\mathbf{2}$ the canonical example against
which both axioms~(10) and~(11) are verified.

% ---------------------------------------------------------
% Definition: Idempotent Element
% ---------------------------------------------------------

\begin{tcolorbox}[colback=propbox, colframe=propborder, arc=2pt,
  left=6pt, right=6pt, top=4pt, bottom=4pt,
  title={\small\textbf{Definition --- Idempotent Element}},
  fonttitle=\small\bfseries]
\begin{definition}[Idempotent Element]
\label{def:idempotent-element}
An element $p$ of a ring $R$ is called \textbf{idempotent} if
\[
p \cdot p = p.
\]
A ring in which every element is idempotent is called an
\textbf{idempotent ring}.
\end{definition}
\end{tcolorbox}

\begin{remark*}[Standard quantified statement]
\[
\begin{aligned}
\operatorname{IdempotentElement}(p, R, \cdot)
&\;\Longleftrightarrow\; p \cdot p = p. \\[4pt]
\operatorname{IdempotentRing}(R, \cdot)
&\;\Longleftrightarrow\; \forall p \in R{:}\; p \cdot p = p.
\end{aligned}
\]
\end{remark*}

\begin{remark*}[Definition predicate reading]
\[
\operatorname{IdempotentElement}(p,R,\cdot)
\qquad
\operatorname{IdempotentRing}(R,\cdot)
\;\Longleftrightarrow\;
\forall p \in R{:}\;
\operatorname{IdempotentElement}(p,R,\cdot)
\]
\end{remark*}

\begin{remark*}[Negated quantified statement]
\[
\neg\,\operatorname{IdempotentElement}(p,R,\cdot)
\;\Longleftrightarrow\;
p \cdot p \neq p.
\]
\end{remark*}

\begin{remark*}[Interpretation]
Idempotence is the single additional axiom that forces Boolean
structure.  In the two-element ring $\mathbf{2}$, verification is
immediate: $0 \cdot 0 = 0$ and $1 \cdot 1 = 1$.  In a power set
algebra $\mathcal{P}(X)$, idempotence of intersection --- $A \cap A
= A$ --- is the set-theoretic counterpart.  The force of the axiom
is not in these examples but in what it implies for an arbitrary
ring: characteristic~2 and commutativity of multiplication both
follow from idempotence alone, as the theorems below establish.
\end{remark*}

% ---------------------------------------------------------
% Definition: Characteristic 2
% ---------------------------------------------------------

\begin{tcolorbox}[colback=propbox, colframe=propborder, arc=2pt,
  left=6pt, right=6pt, top=4pt, bottom=4pt,
  title={\small\textbf{Definition --- Characteristic 2}},
  fonttitle=\small\bfseries]
\begin{definition}[Characteristic 2]
\label{def:characteristic-2}
A ring $R$ is said to have \textbf{characteristic~$2$} if
\[
p + p = 0 \qquad \text{for all } p \in R. \tag{10}
\]
\end{definition}
\end{tcolorbox}

\begin{remark*}[Standard quantified statement]
\[
\operatorname{Characteristic2}(R,+,0)
\;\Longleftrightarrow\;
\forall p \in R{:}\; p + p = 0.
\]
\end{remark*}

\begin{remark*}[Definition predicate reading]
\[
\operatorname{Characteristic2}(R,+,0)
\;\Longleftrightarrow\;
\forall p \in R{:}\;
\operatorname{IdempotentElement}(p, R, +)
\]
\end{remark*}

\begin{remark*}[Negated quantified statement]
\[
\neg\,\operatorname{Characteristic2}(R,+,0)
\;\Longleftrightarrow\;
\exists p \in R{:}\; p + p \neq 0.
\]
\end{remark*}

\begin{remark*}[Interpretation]
Characteristic~2 means that every element is its own additive
inverse: $p + p = 0$ says precisely $p = -p$.  The operation of
negation therefore becomes the identity map --- applying $-$ to
any element returns that element unchanged.  Halmos observes that
in a ring of characteristic~2, the minus sign for additive inverses
is never needed and is henceforth dropped.  Axiom~(7),
$p + (-p) = 0$, is replaced in the official Boolean ring axiom set
by the stronger equation~(10), $p + p = 0$, which is a theorem
derived from idempotence.
\end{remark*}

% ---------------------------------------------------------
% Lemma: Self Additive Inverse (Characteristic 2)
% ---------------------------------------------------------

\begin{lemma}[Self Additive Inverse]
\label{lem:self-additive-inverse}
In any Boolean ring $R$,
\[
p + p = 0 \qquad \text{for all } p \in R.
\]

\smallskip
\noindent
\hyperref[prf:self-additive-inverse]{\textit{Go to proof.}}
\end{lemma}

\begin{remark*}[Standard quantified statement]
\[
\forall p \in R{:}\; p + p = 0.
\]
\end{remark*}

\begin{remark*}[Definition predicate reading]
\[
\operatorname{BooleanRing}(R,+,\cdot,-,0,1)
\;\Longrightarrow\;
\operatorname{Characteristic2}(R,+,0).
\]
\end{remark*}

\begin{remark*}[Interpretation]
This is Halmos consequence~(a): idempotence of multiplication
forces the additive structure to have characteristic~2.  The proof
expands $(p+q)^2$ by distributivity and idempotence, cancels $p$
and $q$, then sets $p = q$ to obtain $p + p = 0$.  The Lean~4
formalisation is \texttt{AdditiveIdempotence} in
\texttt{BooleanAlgebraDefinitions.lean}.
\end{remark*}

% ---------------------------------------------------------
% Lemma: Self Negation
% ---------------------------------------------------------

\begin{lemma}[Self Negation]
\label{lem:self-negation}
In any Boolean ring $R$,
\[
-p = p \qquad \text{for all } p \in R.
\]

\smallskip
\noindent
\hyperref[prf:self-negation]{\textit{Go to proof.}}
\end{lemma}

\begin{remark*}[Standard quantified statement]
\[
\forall p \in R{:}\; -p = p.
\]
\end{remark*}

\begin{remark*}[Interpretation]
This follows immediately from \ref{lem:self-additive-inverse}.
Since $p + p = 0$ and $p + (-p) = 0$ both hold, the uniqueness of
additive inverses gives $-p = p$.  The negation operation is
therefore the identity map on any Boolean ring.  Halmos uses this
to drop the minus sign entirely from his notation for the rest of
the book: since $-p = p$, writing $-p$ and writing $p$ are the
same thing.
\end{remark*}